\chapter{Otimizações de Código Intermediário}

\section{Introdução}

A representação intermediária em árvore (IRT) introduzida no capítulo anterior
não é apenas um formato de saída do \emph{front-end}: ela é o substrato sobre o
qual o compilador aplica transformações que melhoram a qualidade do código sem
alterar seu significado observável.  Essas transformações são denominadas
\textbf{otimizações}.

O termo ``otimização'' é, na prática, um exagero: raramente se garante que o
código produzido é \emph{ótimo} segundo alguma métrica global.  O que se busca
é código \emph{melhor}, ou seja, mais rápido, menor ou mais econômico
em memória,
do que o código gerado diretamente sem nenhuma transformação.

Este capítulo apresenta formalmente dez otimizações aplicadas à IRT
definida no capítulo anterior e, em seguida, descreve a implementação em Haskell
de cada uma delas.

\subsection*{Classificação das Otimizações}

As otimizações são usualmente classificadas pelo escopo de análise:

\begin{description}
  \item[Locais (intra-bloco).] Aplicadas dentro de um único bloco básico, sem
    cruzar desvios condicionais ou rótulos.  São as mais simples e eficientes de
    implementar.  Exemplos: dobramento de constantes, eliminação de
    subexpressões comuns.
  \item[Globais (inter-bloco).] Requerem análise de fluxo de controle entre
    blocos.  Exemplo: propagação de constantes globais, eliminação de código
    morto.
  \item[De laço.] Exploram a estrutura iterativa dos programas.  Exemplos:
    movimentação de código invariante, desenrolamento, fusão de laços.
\end{description}

As sete otimizações cobertas neste capítulo são:

\begin{enumerate}
  \item \textbf{Dobramento de constantes} (\emph{constant folding})
  \item \textbf{Propagação de constantes} (\emph{constant propagation})
  \item \textbf{Dobramento e propagação combinados}
  \item \textbf{Eliminação de subexpressões comuns} (CSE)
  \item \textbf{Eliminação de código morto} (DCE)
  \item \textbf{Movimentação de código invariante de laço} (LICM)
  \item \textbf{Desenrolamento de laço} (\emph{loop unrolling})
  \item \textbf{Fusão de laços} (\emph{loop fusion})
  \item \textbf{Inlining de funções} (\emph{function inlining})
  \item \textbf{Eliminação de chamadas em cauda} (\emph{tail-call elimination})
\end{enumerate}

\subsection*{Terminologia}

Ao longo deste capítulo usamos a seguinte terminologia:

\begin{itemize}
  \item Uma expressão é \textbf{pura} se sua avaliação não produz efeitos
    colaterais: $e$ é pura se $e \in \{\mathbf{CONST}, \mathbf{TEMP},
    \mathbf{NAME}\}$ ou se $e = \mathbf{BINOP}(\mathit{op}, e_1, e_2)$ com
    $e_1$ e $e_2$ puras.  $\mathbf{MEM}$, $\mathbf{CALL}$ e $\mathbf{ESEQ}$
    não são puras.
  \item $\mathit{freeTemps}(e)$ denota o conjunto de temporários que ocorrem
    livres em $e$.
  \item $\mathit{defs}(s)$ denota o conjunto de temporários \emph{definidos}
    (escritos) pelo comando $s$.
  \item $\mathit{uses}(s)$ denota o conjunto de temporários \emph{usados}
    (lidos) pelo comando $s$.
\end{itemize}

% ============================================================================
\section{Dobramento de Constantes}\label{sec:opt-fold}
% ============================================================================

\subsection{Definição Formal}

O \textbf{dobramento de constantes} (\emph{constant folding}) é a avaliação
em tempo de compilação de operações cujos operandos já são conhecidos como
constantes.

\begin{Definition}[Dobramento de Constantes]\label{def:const-fold}
  A função $\mathit{fold} : \mathit{Expr} \to \mathit{Expr}$ é definida
  indutivamente por:
  \begin{align*}
    \mathit{fold}(\mathbf{CONST}(n)) &= \mathbf{CONST}(n) \\
    \mathit{fold}(\mathbf{TEMP}(t))  &= \mathbf{TEMP}(t) \\
    \mathit{fold}(\mathbf{BINOP}(\mathit{op}, e_1, e_2)) &=
    \begin{cases}
      \mathbf{CONST}(\mathit{op}(n_1, n_2))
      & \text{se } \mathit{fold}(e_1) = \mathbf{CONST}(n_1)
      \text{ e } \mathit{fold}(e_2) = \mathbf{CONST}(n_2) \\
      & \quad \text{e } (\mathit{op} \notin \{\mathtt{DIV},\mathtt{MOD}\}
      \text{ ou } n_2 \neq 0) \\
      \mathit{simplify}(\mathit{op}, e_1', e_2')
      & \text{caso contrário}
    \end{cases}
  \end{align*}
  onde $e_i' = \mathit{fold}(e_i)$.  A função $\mathit{simplify}$ aplica
  identidades algébricas:
  \begin{align*}
    \mathit{simplify}(\mathtt{ADD}, e, \mathbf{CONST}(0))     &= e \\
    \mathit{simplify}(\mathtt{ADD}, \mathbf{CONST}(0), e)     &= e \\
    \mathit{simplify}(\mathtt{SUB}, e, \mathbf{CONST}(0))     &= e \\
    \mathit{simplify}(\mathtt{MUL}, \_, \mathbf{CONST}(0))    &=
    \mathbf{CONST}(0) \\
    \mathit{simplify}(\mathtt{MUL}, e, \mathbf{CONST}(1))     &= e \\
    \mathit{simplify}(\mathtt{DIV}, e, \mathbf{CONST}(1))     &= e \\
    \mathit{simplify}(\mathtt{AND}, \_, \mathbf{CONST}(0))    &=
    \mathbf{CONST}(0) \\
    \mathit{simplify}(\mathtt{OR},  e, \mathbf{CONST}(0))     &= e \\
    \mathit{simplify}(\mathit{op}, e_1, e_2)                  &=
    \mathbf{BINOP}(\mathit{op}, e_1, e_2)
  \end{align*}
\end{Definition}

\begin{remark}
  A guarda $n_2 \neq 0$ para $\mathtt{DIV}$ e $\mathtt{MOD}$ evita avaliação
  de divisão por zero em tempo de compilação, preservando o comportamento de
  erro do programa em tempo de execução.
\end{remark}

\begin{example}
  A expressão $\mathbf{BINOP}(\mathtt{ADD},\; \mathbf{CONST}(3),\;
  \mathbf{BINOP}(\mathtt{MUL},\; \mathbf{CONST}(2),\; \mathbf{CONST}(5)))$
  é dobrada em:
  \[
    \mathit{fold}(\cdots) \;=\; \mathbf{CONST}(13)
  \]
  E $\mathbf{BINOP}(\mathtt{MUL},\; \mathbf{TEMP}(x),\; \mathbf{CONST}(0))$
  é simplificada para $\mathbf{CONST}(0)$ pela identidade algébrica, sem
  necessidade de conhecer o valor de $x$.
\end{example}

\subsection{Correção}

A correção do dobramento é garantida pelo seguinte critério: para todo estado
$\sigma$ e expressão $e$,
\[
  \sigma \vdash e \Downarrow v \;\iff\; \sigma \vdash
  \mathit{fold}(e) \Downarrow v.
\]
A demonstração é uma indução estrutural sobre $e$: os casos base
($\mathbf{CONST}$,
$\mathbf{TEMP}$) são triviais; o caso $\mathbf{BINOP}$ segue da correção de
$\mathit{op}(n_1, n_2)$ e das identidades algébricas (que valem em
$\mathbb{Z}$).

% ============================================================================
\section{Propagação de Constantes}\label{sec:opt-prop}
% ============================================================================

\subsection{Definição Formal}

O \textbf{dobramento e propagação combinados} realiza, em uma única passagem
ascendente, tanto a substituição de temporários com valores constantemente
conhecidos quanto o dobramento das subexpressões resultantes.

\begin{Definition}[Ambiente de Propagação]\label{def:prop-env}
  Um \textbf{ambiente de propagação} é uma função parcial
  $\rho : \mathit{Temp} \rightharpoonup \mathbb{Z}$
  que associa a cada temporário (possivelmente) seu valor constante corrente.
\end{Definition}

\begin{Definition}[Propagação Combinada de Expressões]\label{def:fold-prop-expr}
  A função $\mathit{foldProp}_\rho : \mathit{Expr} \to \mathit{Expr}$, para um
  ambiente $\rho$, é definida por:
  \begin{align*}
    \mathit{foldProp}_\rho(\mathbf{TEMP}(t)) &=
    \begin{cases}
      \mathbf{CONST}(\rho(t)) & \text{se } t \in \mathit{dom}(\rho) \\
      \mathbf{TEMP}(t)        & \text{caso contrário}
    \end{cases} \\[4pt]
    \mathit{foldProp}_\rho(\mathbf{BINOP}(\mathit{op}, e_1, e_2)) &=
    \mathit{simplify}\!\left(\mathit{op},\; e_1',\; e_2'\right)
    \quad \text{onde } e_i' = \mathit{foldProp}_\rho(e_i)
  \end{align*}
  (com dobramento imediato quando $e_1', e_2'$ são ambas $\mathbf{CONST}$,
  como na Definição~\ref{def:const-fold}).
\end{Definition}

\begin{Definition}[Propagação de Comandos]\label{def:fold-prop-stmt}
  A função $\mathit{foldPropStmt}$ transforma um comando e atualiza o ambiente:
  \begin{align*}
    \mathit{foldPropStmt}_\rho(\mathbf{SEQ}(s_1, s_2))
    &= \mathbf{SEQ}(s_1',\, s_2'),\; \rho''
    \\
    &\quad \text{onde } (s_1',\, \rho') = \mathit{foldPropStmt}_\rho(s_1)
    \\
    &\quad\quad\quad\quad\; (s_2',\, \rho'') =
    \mathit{foldPropStmt}_{\rho'}(s_2)
    \\[6pt]
    \mathit{foldPropStmt}_\rho(\mathbf{MOVE}(\mathbf{TEMP}(t), e))
    &= \mathbf{MOVE}(\mathbf{TEMP}(t), e'),\; \rho'
    \\
    &\quad \text{onde } e' = \mathit{foldProp}_\rho(e)
    \\
    &\quad\quad\quad\quad\; \rho' =
    \begin{cases}
      \rho[t \mapsto n] & \text{se } e' = \mathbf{CONST}(n) \\
      \rho \setminus \{t\} & \text{caso contrário}
    \end{cases}
  \end{align*}
  Para todos os outros comandos, os ambientes são propagados
  inalterados: $\rho' = \rho$.
\end{Definition}

\begin{example}
  Consideremos a sequência:
  \begin{minted}{text}
  MOVE(TEMP(x), CONST(5))
  MOVE(TEMP(y), BINOP(MUL, TEMP(x), CONST(2)))
  RETURN(TEMP(y))
  \end{minted}
  Com $\rho = \emptyset$ inicialmente:
  \begin{enumerate}
    \item Após o primeiro \textbf{MOVE}: $\rho = \{x \mapsto 5\}$.
    \item O segundo \textbf{MOVE} processa
      $\mathbf{BINOP}(\mathtt{MUL}, \mathbf{TEMP}(x), \mathbf{CONST}(2))$;
      $\mathbf{TEMP}(x)$ é substituído por $\mathbf{CONST}(5)$, e a
      multiplicação
      é dobrada para $\mathbf{CONST}(10)$. $\rho = \{x \mapsto 5, y
      \mapsto 10\}$.
    \item O \textbf{RETURN} propaga: $\mathbf{TEMP}(y) \to \mathbf{CONST}(10)$.
  \end{enumerate}
  O resultado é:
  \begin{minted}{text}
  MOVE(TEMP(x), CONST(5))
  MOVE(TEMP(y), CONST(10))
  RETURN(CONST(10))
  \end{minted}
\end{example}

% ============================================================================
\section{Eliminação de Subexpressões Comuns}\label{sec:opt-cse}
% ============================================================================

\subsection{Definição Formal}

A \textbf{eliminação de subexpressões comuns} (CSE, \emph{Common Subexpression
Elimination}) evita o recomputo de uma expressão pura cujo valor já está
disponível em um temporário.

\begin{Definition}[Tabela CSE]\label{def:cse-table}
  Uma \textbf{tabela CSE} é uma função parcial
  $\theta : \mathit{PureExpr} \rightharpoonup \mathit{Temp}$
  que mapeia cada expressão pura não-trivial ao temporário que contém seu valor
  mais recentemente calculado.
  A tabela é válida dentro de um \emph{bloco básico}: ela é reiniciada a cada
  fronteira de bloco ($\mathbf{LABEL}$, $\mathbf{JUMP}$, $\mathbf{CJUMP}$,
  $\mathbf{RETURN}$).
\end{Definition}

\begin{Definition}[Transformação CSE]\label{def:cse-transform}
  Para cada instrução $\mathbf{MOVE}(\mathbf{TEMP}(t), e)$ onde $e$ é pura e
  não-trivial:
  \[
    \mathit{CSE}_\theta(\mathbf{MOVE}(\mathbf{TEMP}(t), e)) =
    \begin{cases}
      \mathbf{MOVE}(\mathbf{TEMP}(t), \mathbf{TEMP}(\theta(e)))
      & \text{se } e \in \mathit{dom}(\theta) \\[4pt]
      \begin{array}{l}
        \mathbf{MOVE}(\mathbf{TEMP}(c), e) \\
        \mathbf{MOVE}(\mathbf{TEMP}(t), \mathbf{TEMP}(c))
      \end{array}
      & \text{caso contrário}
    \end{cases}
  \]
  No segundo caso, $c$ é um novo temporário introduzido pelo compilador, e
  $\theta$ é atualizado com $\theta' = \theta[e \mapsto c]$.
  Em ambos os casos, todas as entradas de $\theta$ que mencionam $t$ são
  removidas (pois $t$ acaba de ser redefinido).
\end{Definition}

\begin{remark}
  A restrição a expressões \emph{puras} é essencial: $\mathbf{MEM}(e)$ pode
  produzir valores diferentes após uma escrita em memória, e $\mathbf{CALL}$
  pode ter efeitos colaterais.  A restrição ao escopo \emph{intra-bloco}
  evita problemas com caminhos alternativos que podem não executar o cálculo
  original.
\end{remark}

\begin{example}
  A sequência:
  \begin{minted}{text}
  MOVE(TEMP(a), BINOP(ADD, TEMP(x), TEMP(y)))
  MOVE(TEMP(b), BINOP(ADD, TEMP(x), TEMP(y)))
  \end{minted}
  é transformada em:
  \begin{minted}{text}
  MOVE(TEMP(_cse0), BINOP(ADD, TEMP(x), TEMP(y)))
  MOVE(TEMP(a),     TEMP(_cse0))
  MOVE(TEMP(b),     TEMP(_cse0))
  \end{minted}
  A expressão $\mathbf{BINOP}(\mathtt{ADD}, x, y)$ é calculada uma única vez.
\end{example}

% ============================================================================
\section{Eliminação de Código Morto}\label{sec:opt-dce}
% ============================================================================

\subsection{Definição Formal}

A \textbf{eliminação de código morto} (DCE, \emph{Dead Code Elimination})
remove instruções que nunca influenciam o resultado observável do programa.
Implementamos duas fases complementares.

\subsubsection*{Fase 1: Eliminação de Código Inalcançável}

\begin{Definition}[Código Inalcançável]\label{def:unreachable}
  Um comando $s$ é \textbf{inalcançável} se toda execução que o antecede passou
  por um $\mathbf{JUMP}$ ou $\mathbf{RETURN}$ sem um rótulo intermediário.
  Formalmente, uma instrução na posição $i$ da sequência linearizada é
  inalcançável se a posição $i-1$ é ocupada por $\mathbf{JUMP}$ ou
  $\mathbf{RETURN}$ e a posição $i$ não é $\mathbf{LABEL}$.
\end{Definition}

A fase 1 percorre a sequência linearizada e elimina todos os comandos
inalcançáveis, preservando os $\mathbf{LABEL}$s que podem ser destinos de
saltos provenientes de outros pontos do programa.

\subsubsection*{Fase 2: Eliminação de Atribuições Mortas}

\begin{Definition}[Atribuição Morta]\label{def:dead-assign}
  Uma instrução $\mathbf{MOVE}(\mathbf{TEMP}(t), e)$ onde $e$ é pura é
  \textbf{morta} se $t$ é redefinido antes de ser lido, ou se a função
  termina (via \textbf{RETURN}) sem que $t$ seja lido.
\end{Definition}

A fase 2 realiza uma análise \emph{forward} com uma tabela de
\emph{definições pendentes}:

\begin{Definition}[Estado DAE]\label{def:dae-state}
  O estado da análise é um par $(\mathit{pend}, \mathit{dead})$ onde:
  \begin{itemize}
    \item $\mathit{pend} : \mathit{Temp} \rightharpoonup \mathit{Idx}
      \times \mathit{Expr}$ mapeia cada temporário ao índice e expressão
      de sua atribuição mais recente ainda não consumida;
    \item $\mathit{dead} \subseteq \mathit{Idx}$ é o conjunto de índices das
      instruções marcadas para remoção.
  \end{itemize}
\end{Definition}

As regras de atualização do estado são:

\begin{align*}
  \text{(novo MOVE)} &\quad
  \mathbf{MOVE}(\mathbf{TEMP}(t), e) \text{ no índice } i: \\
  &\quad\quad \text{consumir } \mathit{freeTemps}(e) \text{ de }
  \mathit{pend}; \\
  &\quad\quad \text{se } t \in \mathit{pend} \text{ com expr. pura} \to
  \mathit{dead} \mathrel{+}= \mathit{pend}(t).\mathit{idx}; \\
  &\quad\quad \mathit{pend}(t) \leftarrow (i, e) \\[4pt]
  \text{(RETURN)} &\quad
  \text{consumir usos; todos os } t \in \mathit{pend} \text{ com expr.
  pura} \to \mathit{dead}; \\
  &\quad\quad \mathit{pend} \leftarrow \emptyset \\[4pt]
  \text{(LABEL, JUMP, CJUMP)} &\quad
  \text{consumir usos; } \mathit{pend} \leftarrow \emptyset
  \quad \text{(descarte conservador)}
\end{align*}

\begin{remark}
  O descarte conservador em fronteiras de bloco é necessário porque, sem
  análise de fluxo inter-bloco, não sabemos quais temporários são lidos em
  destinos de salto.  Atribuições impuras (com $\mathbf{CALL}$, $\mathbf{MEM}$)
  nunca são marcadas mortas, preservando seus efeitos colaterais.
\end{remark}

\begin{example}
  Após a CSE do exemplo anterior, pode restar:
  \begin{minted}{text}
  MOVE(TEMP(x), CONST(5))
  MOVE(TEMP(y), CONST(10))
  RETURN(CONST(10))
  \end{minted}
  A DCE detecta que \texttt{x} e \texttt{y} são definidos mas nunca lidos
  antes do \textbf{RETURN} (que usa diretamente $\mathbf{CONST}(10)$),
  eliminando as duas atribuições:
  \begin{minted}{text}
  RETURN(CONST(10))
  \end{minted}
\end{example}

% ============================================================================
\section{Movimentação de Código Invariante de Laço}\label{sec:opt-licm}
% ============================================================================

\subsection{Estrutura de Laço na IRT}

A IRT linearizada representa um laço \texttt{while} pelo padrão de seis
elementos consecutivos:

\[
  \mathbf{LABEL}(\ell_h) \;\;
  \mathbf{CJUMP}(e,\, \ell_b,\, \ell_f) \;\;
  \mathbf{LABEL}(\ell_b) \;\;
  \underbrace{s_1\; \cdots\; s_k}_{\text{corpo}} \;\;
  \mathbf{JUMP}(\mathbf{NAME}(\ell_h)) \;\;
  \mathbf{LABEL}(\ell_f)
\]

onde $\ell_h$ é a cabeça do laço, $\ell_b$ é o início do corpo, $\ell_f$ é o
ponto de saída, e a aresta de retorno $\mathbf{JUMP}(\mathbf{NAME}(\ell_h))$
fecha o laço.

\begin{Definition}[Representação de Laço]\label{def:loop}
  Um \textbf{laço} $L$ é uma quíntupla
  $(\ell_h, \ell_b, \ell_f, e_{\mathit{cond}}, B)$ onde $B = [s_1, \ldots, s_k]$
  é a lista de instruções do corpo.  Definimos:
  \begin{align*}
    \mathit{defs}(L) &= \{ t \mid \mathbf{MOVE}(\mathbf{TEMP}(t), \_)
    \in B \} \\
    \mathit{condVars}(L) &= \mathit{freeTemps}(e_{\mathit{cond}})
  \end{align*}
\end{Definition}

\subsection{Definição Formal de LICM}

\begin{Definition}[Código Invariante de Laço]\label{def:licm}
  Uma instrução $\mathbf{MOVE}(\mathbf{TEMP}(t), e)$ na posição $i$ de $B$
  é \textbf{invariante de laço} se:
  \begin{enumerate}
    \item $e$ é pura;
    \item $\mathit{freeTemps}(e) \cap \mathit{defs}(L) = \emptyset$
      (nenhum temporário de $e$ é modificado no laço);
    \item $t$ é definido exatamente uma vez em $B$;
    \item $t \notin \mathit{uses}(s_1) \cup \cdots \cup \mathit{uses}(s_{i-1})$
      (nenhuma instrução anterior ao ponto de definição lê $t$ no corpo).
  \end{enumerate}
\end{Definition}

A condição (3) garante que mover $t$ para antes do laço não altera o valor
visto em iterações subsequentes; a condição (4) garante que a primeira iteração
não lê $t$ antes de sua definição original.

A transformação LICM substitui o laço original por:

\[
  \underbrace{s_{i_1} \;\cdots\; s_{i_m}}_{\text{código içado}} \;\;
  \mathbf{LABEL}(\ell_h) \;\; \mathbf{CJUMP}(e, \ell_b, \ell_f) \;\;
  \mathbf{LABEL}(\ell_b) \;\;
  \underbrace{B \setminus \{s_{i_1},\ldots,s_{i_m}\}}_{\text{corpo
  reduzido}} \;\;
  \mathbf{JUMP}(\mathbf{NAME}(\ell_h)) \;\; \mathbf{LABEL}(\ell_f)
\]

onde $s_{i_1},\ldots,s_{i_m}$ são as instruções invariantes, na sua ordem de
aparecimento em $B$.

\begin{example}
  O corpo do laço:
  \begin{minted}{text}
  MOVE(TEMP(k), BINOP(MUL, CONST(4), CONST(8)))
  MOVE(TEMP(a), BINOP(MUL, TEMP(i), TEMP(k)))
  MOVE(TEMP(i), BINOP(SUB, TEMP(i), CONST(1)))
  \end{minted}
  A instrução \texttt{MOVE(TEMP(k), ...)} é invariante: sua expressão é pura,
  não usa nenhum temporário modificado no laço, e \texttt{k} é definido
  exatamente uma vez.  Após LICM:
  \begin{minted}{text}
  MOVE(TEMP(k), BINOP(MUL, CONST(4), CONST(8)))   -- içado
  LABEL(loop_head)
  CJUMP(..., loop_body, loop_exit)
  LABEL(loop_body)
    MOVE(TEMP(a), BINOP(MUL, TEMP(i), TEMP(k)))
    MOVE(TEMP(i), BINOP(SUB, TEMP(i), CONST(1)))
  JUMP(NAME(loop_head))
  LABEL(loop_exit)
  \end{minted}
  Se combinada com dobramento de constantes, \texttt{k} seria eliminado e o
  literal $32$ seria usado diretamente.
\end{example}

% ============================================================================
\section{Desenrolamento de Laço}\label{sec:opt-unroll}
% ============================================================================

\subsection{Definição Formal}

O \textbf{desenrolamento de laço} (\emph{loop unrolling}) com fator $k$ replica
o corpo do laço $k$ vezes dentro de uma única iteração, reduzindo o número de
avaliações da condição e abrindo oportunidades para outras otimizações.

\begin{Definition}[Desenrolamento com Fator $k$]\label{def:unroll}
  Seja $L = (\ell_h, \ell_b, \ell_f, e, B)$ um laço.
  Para $k \geq 2$, o desenrolamento de $L$ com fator $k$ produz:
  \[
    \mathbf{LABEL}(\ell_h)\;\;
    \mathbf{CJUMP}(e,\, \ell_b,\, \ell_f)\;\;
    \mathbf{LABEL}(\ell_b)\;\;
    B\;\;
    C_1 \;\; C_2 \;\; \cdots \;\; C_{k-1} \;\;
    \mathbf{JUMP}(\mathbf{NAME}(\ell_h))\;\;
    \mathbf{LABEL}(\ell_f)
  \]
  onde cada cópia extra $C_j$ ($1 \leq j \leq k-1$) é:
  \[
    C_j = \mathbf{CJUMP}(e,\, \ell_h^{(j)},\, \ell_f)\;\;
    \mathbf{LABEL}(\ell_h^{(j)})\;\;
    B[\ell \mapsto \ell\text{\_u}j \mid \ell \in \mathit{InternalLabels}(B)]
  \]
  $\ell_h^{(j)}$ é o rótulo de entrada da $j$-ésima cópia extra, e a renomeação
  de rótulos internos ($\ell \mapsto \ell\text{\_u}j$) evita colisões entre
  cópias.
\end{Definition}

\begin{remark}
  Cada cópia extra reavalia a condição $e$ antes de executar o corpo copiado.
  Isso garante que, se o laço terminar numa iteração que não seja múltiplo de
  $k$, a cópia extra detecta a saída e salta para $\ell_f$.  O corpo original
  (cópia 0) permanece sob o único $\mathbf{CJUMP}$ do cabeçalho.
\end{remark}

\begin{example}
  Para o laço com corpo $[s_1, s_2]$ e fator $k = 2$, o resultado é:
  \begin{minted}{text}
  LABEL(lh)
  CJUMP(e, lb, lf)
  LABEL(lb)
    s1 ; s2                              -- cópia 0
    CJUMP(e, lh_unroll_1, lf)           -- guarda extra
    LABEL(lh_unroll_1)
    s1_u1 ; s2_u1                       -- cópia 1 (rótulos renomeados)
  JUMP(NAME(lh))
  LABEL(lf)
  \end{minted}
\end{example}

% ============================================================================
\section{Fusão de Laços}\label{sec:opt-fusion}
% ============================================================================

\subsection{Definição Formal}

A \textbf{fusão de laços} (\emph{loop fusion}) combina dois laços adjacentes
com a mesma condição de iteração em um único laço, reduzindo a sobrecarga de
controle e melhorando a localidade de cache.

\begin{Definition}[Adjacência]\label{def:adjacent}
  Dois laços $L_1$ e $L_2$ são \textbf{adjacentes} se o índice do
  $\mathbf{LABEL}(\ell_f^{L_1})$ (saída de $L_1$) é imediatamente seguido pelo
  índice do $\mathbf{LABEL}(\ell_h^{L_2})$ (cabeça de $L_2$):
  \[
    \mathit{exitIdx}(L_1) + 1 = \mathit{headIdx}(L_2).
  \]
\end{Definition}

\begin{Definition}[Fusibilidade]\label{def:fusible}
  Dois laços adjacentes $L_1 = (\ell_h^1, \ell_b^1, \ell_f^1, e_1, B_1)$ e
  $L_2 = (\ell_h^2, \ell_b^2, \ell_f^2, e_2, B_2)$ são \textbf{fusíveis} se:
  \begin{enumerate}
    \item $e_1 = e_2$ (mesma condição de iteração, igualdade estrutural);
    \item $\mathit{last}(B_1) = \mathit{last}(B_2)$ (mesma instrução de
      atualização da variável de indução no final de cada corpo);
    \item \textbf{Ausência de dependência cruzada:}
      \[
        \bigl(\mathit{defs}(L_1) \setminus \mathit{condVars}(L_1)\bigr)
        \cap \mathit{firstUses}(B_2) = \emptyset
      \]
      onde $\mathit{firstUses}(B_2)$ é o conjunto de temporários lidos em $B_2$
      antes de serem escritos (definição abaixo).
  \end{enumerate}
\end{Definition}

\begin{Definition}[Primeiros Usos]\label{def:first-uses}
  $\mathit{firstUses}(B)$ é computado por um acúmulo forward sobre $B$:
  partindo de $D = \emptyset$, $U = \emptyset$, para cada instrução $s \in B$:
  \[
    U \leftarrow U \cup (\mathit{uses}(s) \setminus D),
    \qquad
    D \leftarrow D \cup \mathit{defs}(s).
  \]
  O resultado é $U$ ao final da lista.
\end{Definition}

Quando $L_1$ e $L_2$ são adjacentes e fusíveis, o laço fundido é:

\[
  \mathbf{LABEL}(\ell_h^1)\;\;
  \mathbf{CJUMP}(e_1,\, \ell_b^1,\, \ell_f^1)\;\;
  \mathbf{LABEL}(\ell_b^1)\;\;
  \underbrace{\mathit{init}(B_1)}_{\text{corpo de }L_1\text{ sem a
  última instrução}} \;\;
  B_2 \;\;
  \mathbf{JUMP}(\mathbf{NAME}(\ell_h^1))\;\;
  \mathbf{LABEL}(\ell_f^1)
\]

A última instrução de $B_1$ (atualização da variável de indução) é removida de
$B_1$ porque ela já aparece no final de $B_2$.  Isso garante que a atualização
da variável de indução ocorre exatamente uma vez por iteração fundida.

\begin{example}
  Dois laços adjacentes que acumulam duas somas independentes:
  \begin{minted}{text}
  -- Laço 1
  LABEL(lh1) ; CJUMP(BINOP(GT, TEMP(n), CONST(0)), lb1, lf1)
  LABEL(lb1)
    MOVE(TEMP(s1), BINOP(ADD, TEMP(s1), TEMP(a)))
    MOVE(TEMP(n),  BINOP(SUB, TEMP(n),  CONST(1)))
  JUMP(NAME(lh1)) ; LABEL(lf1)

  -- Laço 2
  LABEL(lh2) ; CJUMP(BINOP(GT, TEMP(n), CONST(0)), lb2, lf2)
  LABEL(lb2)
    MOVE(TEMP(s2), BINOP(ADD, TEMP(s2), TEMP(b)))
    MOVE(TEMP(n),  BINOP(SUB, TEMP(n),  CONST(1)))
  JUMP(NAME(lh2)) ; LABEL(lf2)
  \end{minted}
  Os dois laços possuem a mesma condição e o mesmo passo.  A verificação de
  dependência cruzada: $\mathit{defs}(L_1) \setminus \{n\} = \{s1\}$, e
  $s1 \notin \mathit{firstUses}(B_2)$.  Portanto são fusíveis:
  \begin{minted}{text}
  LABEL(lh1) ; CJUMP(BINOP(GT, TEMP(n), CONST(0)), lb1, lf1)
  LABEL(lb1)
    MOVE(TEMP(s1), BINOP(ADD, TEMP(s1), TEMP(a)))
    -- init(B1): sem o MOVE(n, ...)
    MOVE(TEMP(s2), BINOP(ADD, TEMP(s2), TEMP(b)))
    MOVE(TEMP(n),  BINOP(SUB, TEMP(n),  CONST(1)))
  JUMP(NAME(lh1)) ; LABEL(lf1)
  \end{minted}
\end{example}

% ============================================================================
\section{Inlining de Funções}\label{sec:opt-inline}
% ============================================================================

\subsection{Definição Formal}

O \textbf{inlining de funções} substitui uma chamada a uma função pelo corpo
desta, eliminando a sobrecarga da chamada e expondo o código do sítio de chamada
a outras otimizações (dobramento, propagação, CSE).

\begin{Definition}[Substituição de Inlining]\label{def:inline}
  Dada uma função $f$ com parâmetros formais $\bar{x} = (x_1, \ldots, x_n)$ e
  corpo $B_f$, e uma chamada $\mathbf{CALL}(\mathbf{NAME}(f), [a_1,\ldots,a_n])$
  num sítio de chamada com sufixo de frescura $s$, a \textbf{substituição de
  inlining} produz:
  \[
    \begin{aligned}
      \mathit{inline}(f, [a_1,\ldots,a_n])
      &= \mathbf{ESEQ}( \\
        &\quad \mathbf{MOVE}(\mathbf{TEMP}(x_1^s), a_1) \mathbin{;}
        \cdots \mathbin{;}
        \mathbf{MOVE}(\mathbf{TEMP}(x_n^s), a_n) \\
        &\quad \mathbin{;} B_f^s \mathbin{;}
        \mathbf{LABEL}(\ell_\mathit{exit}^s), \\
      &\quad \mathbf{TEMP}(r^s))
    \end{aligned}
  \]
  onde:
  \begin{itemize}
    \item $x_i^s = x_i \mathbin{++} s$: parâmetros renomeados com frescura;
    \item $B_f^s$: corpo de $f$ com todos os temporários e rótulos locais
      renomeados com o sufixo $s$;
    \item cada $\mathbf{RETURN}([e])$ em $B_f^s$ é substituído por
      $\mathbf{MOVE}(\mathbf{TEMP}(r^s), e) \mathbin{;}
      \mathbf{JUMP}(\mathbf{NAME}(\ell_\mathit{exit}^s))$;
    \item $r^s$ e $\ell_\mathit{exit}^s$ são temporário e rótulo novos para
      capturar o valor de retorno e reunificar o fluxo de controle.
  \end{itemize}
  Para chamadas \emph{void} (em posição de comando), o envoltório
  $\mathbf{ESEQ}(\cdot, \mathbf{TEMP}(r^s))$ é suprimido.
\end{Definition}

A renomeação de locais é essencial: a mesma função pode ser inlinada
em múltiplos
sítios, e os sufixos distintos evitam colisões de temporários entre cópias.

\begin{Definition}[Elegibilidade para Inlining]\label{def:inline-eligible}
  Uma chamada $\mathbf{CALL}(\mathbf{NAME}(f), \bar{a})$ é \textbf{elegível
  para inlining} se:
  \begin{enumerate}
    \item O alvo é estático: $\mathbf{NAME}(f)$, não uma expressão computada;
    \item $f$ não é auto-recursiva: $f \notin \mathit{calls}(B_f)$, onde
      $\mathit{calls}$ coleta todos os $\mathbf{NAME}$ em posição de chamada;
    \item O tamanho do corpo não excede um limiar $\tau$:
      $|\mathit{linearize}(B_f)| \leq \tau$.
  \end{enumerate}
\end{Definition}

A condição (2) impede que o inlining de funções recursivas gere
expansão infinita.
A condição (3) é uma heurística: funções grandes não devem ser inlinadas porque
o aumento de código (\emph{code bloat}) pode degradar a localidade de cache e
aumentar o tamanho do binário.

\begin{Theorem}[Correção do Inlining]\label{thm:inline-correct}
  Para toda função $f$ elegível e chamada
  $\mathbf{CALL}(\mathbf{NAME}(f), [\bar{a}])$:
  \[
    \sigma \vdash \mathbf{CALL}(\mathbf{NAME}(f), [\bar{a}]) \Downarrow v
    \;\iff\;
    \sigma \vdash \mathit{inline}(f, [\bar{a}]) \Downarrow v.
  \]
\end{Theorem}

A correção segue da regra semântica de invocação: avaliar os argumentos, ligar
aos parâmetros e executar o corpo é exatamente o que $\mathit{inline}$ realiza
estaticamente.

Após inlinar todas as chamadas elegíveis, funções que deixam de ser
referenciadas
são removidas (\emph{eliminação de funções mortas}).

% ============================================================================
\section{Eliminação de Chamadas em Cauda}\label{sec:opt-tco}
% ============================================================================

\subsection{Chamadas em Posição de Cauda}

Uma chamada está em \textbf{posição de cauda} numa função $g$ se o único efeito
observável após sua execução é retornar seu resultado ao chamador de $g$.  Nesse
caso, o frame de ativação corrente pode ser \emph{reutilizado} em vez
de empilhar
um novo frame, convertendo recursão em iteração com consumo de pilha $O(1)$.

\begin{Definition}[Chamada em Cauda na IRT]\label{def:tail-call}
  Na sequência linearizada de uma função, dois padrões constituem chamadas em
  cauda:
  \begin{align*}
    \text{(com valor de retorno)} &\quad
    \mathbf{RETURN}([\mathbf{CALL}(\mathbf{NAME}(f), \bar{a})]) \\
    \text{(sem valor de retorno)} &\quad
    \mathbf{EXP}(\mathbf{CALL}(\mathbf{NAME}(f), \bar{a})) \mathbin{;}
    \mathbf{RETURN}([])
  \end{align*}
\end{Definition}

\subsection{Eliminação de Chamada Recursiva em Cauda}

Quando $f = g$ (auto-recursão em cauda), o frame de $g$ pode ser completamente
reutilizado: atualiza-se os parâmetros e salta-se para o início do corpo.

\begin{Definition}[Transformação TCO]\label{def:tco}
  Seja $g$ uma função com parâmetros $\bar{x} = (x_1, \ldots, x_n)$,
  corpo linearizado $S$, e $\ell_h = \mathtt{L\_tco\_}g$ um novo rótulo.
  A transformação TCO de $g$ produz:
  \begin{enumerate}
    \item Inserir $\mathbf{LABEL}(\ell_h)$ no início de $S$.
    \item Substituir cada padrão de cauda (Definição~\ref{def:tail-call}) por:
      \begin{align*}
        &\mathbf{MOVE}(\mathbf{TEMP}(t_1), a_1) \mathbin{;} \cdots \mathbin{;}
        \mathbf{MOVE}(\mathbf{TEMP}(t_n), a_n) & \text{(passo 1:
        avaliar argumentos)} \\
        &\mathbf{MOVE}(\mathbf{TEMP}(x_1), \mathbf{TEMP}(t_1))
        \mathbin{;} \cdots \mathbin{;}
        \mathbf{MOVE}(\mathbf{TEMP}(x_n), \mathbf{TEMP}(t_n)) &
        \text{(passo 2: atualizar parâmetros)} \\
        &\mathbf{JUMP}(\mathbf{NAME}(\ell_h)) & \text{(passo 3:
        reiniciar o laço)}
      \end{align*}
      onde $t_1, \ldots, t_n$ são novos temporários auxiliares.
  \end{enumerate}
\end{Definition}

O \textbf{passo duplo} (avaliar para auxiliares, depois copiar para parâmetros)
é necessário para evitar o \emph{hazard} de leitura após escrita.  Por exemplo,
se $\bar{a} = (x_2, x_1)$ e executarmos diretamente
$\mathbf{MOVE}(\mathbf{TEMP}(x_1), x_2) \mathbin{;}
\mathbf{MOVE}(\mathbf{TEMP}(x_2), x_1)$,
a segunda instrução leria o valor \emph{já sobrescrito} de $x_1$.  Os
temporários
auxiliares (gerados como $\mathtt{\_tco\_0}, \ldots, \mathtt{\_tco\_}(n{-}1)$)
quebram essa dependência cíclica.

\begin{example}
  Considere uma função fatorial com acumulador:
  \begin{minted}{text}
  func factAcc(n, acc):
    ...
    RETURN [CALL(NAME(factAcc), [BINOP(BSub, TEMP(n), CONST(1)),
                                 BINOP(BMul, TEMP(acc), TEMP(n))])]
  \end{minted}
  Após TCO:
  \begin{minted}{text}
  func factAcc(n, acc):
    LABEL(L_tco_factAcc)
    ...
    MOVE(TEMP(_tco_0), BINOP(BSub, TEMP(n),   CONST(1)))
    MOVE(TEMP(_tco_1), BINOP(BMul, TEMP(acc), TEMP(n)))
    MOVE(TEMP(n),   TEMP(_tco_0))
    MOVE(TEMP(acc), TEMP(_tco_1))
    JUMP(NAME(L_tco_factAcc))
  \end{minted}
  A recursão é convertida num laço sem crescimento de pilha.
\end{example}

\begin{Theorem}[Correção da TCO]\label{thm:tco-correct}
  Seja $g$ uma função com auto-recursão em cauda.  Para todo estado inicial
  $\sigma$, $g$ e $\mathit{tco}(g)$ produzem o mesmo valor, ou ambos divergem.
\end{Theorem}

A demonstração usa indução sobre o número de chamadas recursivas: no caso base
o laço termina na primeira iteração; no passo indutivo, cada volta ao
$\mathbf{LABEL}(\ell_h)$ corresponde exatamente a uma chamada recursiva com
parâmetros atualizados, preservando a semântica de substituição.

% ============================================================================
\section{Implementação em Haskell}\label{sec:opt-impl}
% ============================================================================

Esta seção descreve a implementação concreta das otimizações em Haskell,
organizada nos módulos do pacote \texttt{bcc328}.

\subsection{Infraestrutura Comum:
\texttt{IR.Opt.Loop.Loop}}\label{sec:impl-loop}

O módulo \texttt{IR.Opt.Loop.Loop} define a estrutura de dados \texttt{Loop} e
os utilitários compartilhados pelas três otimizações de laço.

\begin{minted}{haskell}
data Loop = Loop
  { loopHeadLabel :: Label
  , loopBodyLabel :: Label
  , loopExitLabel :: Label
  , loopCond      :: Expr
  , loopBodyStmts :: [Stmt]
  , loopHeadIdx   :: Int
  , loopExitIdx   :: Int
  }
\end{minted}

A função \texttt{findLoops} varre a sequência linearizada procurando o padrão
\texttt{LABEL lh / CJUMP / LABEL lnext / ... / JUMP(NAME lh) / LABEL lf}.
O módulo aceita ambas as convenções de CJUMP (salto verdadeiro para o corpo ou
para a saída), determinando qual é o corpo pelo rótulo que segue imediatamente
o CJUMP:

\begin{minted}{haskell}
go i (LABEL lh : CJUMP cond l1 l2 : LABEL lnext : rest)
  | lnext == l1 || lnext == l2 =
      let (lb, lf) = if lnext == l1 then (l1, l2) else (l2, l1)
      in case splitAtBackEdge lh rest of
           Just (body, LABEL lf' : remaining) | lf' == lf ->
             loop : go (exitIdx + 1) remaining
           _ -> ...
\end{minted}

Os utilitários \texttt{linearize} e \texttt{rebuildSeq} convertem entre a
representação em árvore SEQ e a lista plana, simplificando o processamento:

\begin{minted}{haskell}
linearize :: Stmt -> [Stmt]
linearize (SEQ s1 s2) = linearize s1 ++ linearize s2
linearize s           = [s]

rebuildSeq :: [Stmt] -> Stmt
rebuildSeq []     = EXP (CONST 0)
rebuildSeq [s]    = s
rebuildSeq (s:ss) = SEQ s (rebuildSeq ss)
\end{minted}

\subsection{Dobramento de Constantes:
\texttt{IR.Opt.ConstFold}}\label{sec:impl-fold}

O módulo \texttt{IR.Opt.ConstFold} implementa uma passagem puramente
estrutural, sem estado, que percorre a árvore de baixo para cima:

\begin{minted}{haskell}
foldExpr :: Expr -> Expr
foldExpr (BINOP op e1 e2) =
  let e1' = foldExpr e1
      e2' = foldExpr e2
  in case (e1', e2') of
       (CONST a, CONST b)
         | op `notElem` [BDiv, BMod] || b /= 0
           -> CONST (applyBinOp op a b)
       _ -> simplify op e1' e2'
foldExpr (MEM e)      = MEM (foldExpr e)
foldExpr (CALL ef as) = CALL (foldExpr ef) (map foldExpr as)
foldExpr (ESEQ s e)   = ESEQ (foldStmt s) (foldExpr e)
foldExpr e            = e
\end{minted}

A função \texttt{simplify} é exportada para reutilização pelo módulo
combinado, evitando duplicação de código.  Ela codifica as identidades
algébricas da Definição~\ref{def:const-fold} como equações de padrão Haskell.

\subsection{Propagação Combinada:
\texttt{IR.Opt.ConstFoldProp}}\label{sec:impl-foldprop}

O módulo \texttt{IR.Opt.ConstFoldProp} implementa a passagem combinada com um
ambiente explícito.  O tipo \texttt{Env = Map Temp Int} representa $\rho$:

\begin{minted}{haskell}
foldPropExpr :: Env -> Expr -> Expr
foldPropExpr env (TEMP t) =
  case Map.lookup t env of
    Just n  -> CONST n
    Nothing -> TEMP t
foldPropExpr env (BINOP op e1 e2) =
  let e1' = foldPropExpr env e1
      e2' = foldPropExpr env e2
  in case (e1', e2') of
       (CONST a, CONST b)
         | op `notElem` [BDiv, BMod] || b /= 0
           -> CONST (applyBinOp op a b)
       _ -> simplify op e1' e2'
foldPropExpr env (MEM e)      = MEM (foldPropExpr env e)
foldPropExpr env (CALL ef as) =
  CALL (foldPropExpr env ef) (map (foldPropExpr env) as)
foldPropExpr _   e            = e
\end{minted}

A função \texttt{foldPropStmt} encadeia o ambiente de esquerda para direita
através de sequências, atualizando $\rho$ a cada atribuição a temporário:

\begin{minted}{haskell}
foldPropStmt :: Env -> Stmt -> (Stmt, Env)
foldPropStmt env (SEQ s1 s2) =
  let (s1', env1) = foldPropStmt env  s1
      (s2', env2) = foldPropStmt env1 s2
  in (SEQ s1' s2', env2)
foldPropStmt env (MOVE (TEMP t) src) =
  let src' = foldPropExpr env src
      env'  = case src' of
                CONST n -> Map.insert t n env
                _       -> Map.delete t env
  in (MOVE (TEMP t) src', env')
\end{minted}

\subsection{CSE: \texttt{IR.Opt.CSE}}\label{sec:impl-cse}

O módulo \texttt{IR.Opt.CSE} mantém o estado \texttt{CSEState} com a tabela
$\theta$ e um contador para gerar novos temporários CSE:

\begin{minted}{haskell}
data CSEState = CSEState
  { cseTable   :: Map Expr Temp
  , cseCounter :: Int
  }
\end{minted}

A função central \texttt{cseOneStmt} implementa a
Definição~\ref{def:cse-transform}:

\begin{minted}{haskell}
cseOneStmt st (MOVE (TEMP t) e)
  | isPure e && not (isTrivial e) =
      case Map.lookup e (cseTable st) of
        Just cseT ->
          ([MOVE (TEMP t) (TEMP cseT)], killTemp t st)
        Nothing ->
          let k    = cseCounter st
              cseT = "_cse" ++ show k
              tbl1 = Map.insert e cseT (cseTable st)
              st1  = st { cseCounter = k+1, cseTable = tbl1 }
          in ( [MOVE (TEMP cseT) e, MOVE (TEMP t) (TEMP cseT)]
             , killTemp t st1 )
cseOneStmt st s | isBlockBoundary s =
  ([s], st { cseTable = Map.empty })
cseOneStmt st s = ([s], st)
\end{minted}

A função \texttt{killTemp t} remove da tabela todas as entradas cuja
expressão menciona $t$, garantindo que entradas obsoletas não sejam reutilizadas
após uma redefinição.

\subsection{DCE: \texttt{IR.Opt.DCE}}\label{sec:impl-dce}

O módulo \texttt{IR.Opt.DCE} implementa as duas fases sequencialmente:

\begin{minted}{haskell}
dceFuncDef fd = fd { funcBody = process (funcBody fd) }
  where
    process = rebuildSeq . deadAssignElim
            . removeUnreachable . linearize
\end{minted}

A fase 1 (\texttt{removeUnreachable}) percorre a lista e descarta comandos
entre um $\mathbf{JUMP}$/$\mathbf{RETURN}$ e o próximo $\mathbf{LABEL}$:

\begin{minted}{haskell}
removeUnreachable (s:ss)
  | isHardJump s = s : skipUntilLabel ss
  | otherwise    = s : removeUnreachable ss
\end{minted}

A fase 2 (\texttt{deadAssignElim}) usa \texttt{Data.List.foldl'} para processar
os comandos indexados, acumulando o estado DAE:

\begin{minted}{haskell}
deadAssignElim stmts =
  let indexed    = zip [0..] stmts
      finalState = foldl' step initDAE indexed
      extraDead  = Set.fromList
        [ pendingIdx p | p <- Map.elems (pending finalState)
                       , isPure (pendingExpr p) ]
      allDead    = dead finalState `Set.union` extraDead
  in [ s | (i, s) <- indexed, i `Set.notMember` allDead ]
\end{minted}

\subsection{LICM: \texttt{IR.Opt.Loop.LICM}}\label{sec:impl-licm}

A implementação de LICM itera sobre os laços detectados, aplicando
\texttt{hoistInvariants} a cada um:

\begin{minted}{haskell}
hoistInvariants :: Loop -> ([Stmt], [Stmt])
hoistInvariants loop = go [] [] [] (loopBodyStmts loop)
  where
    defs = loopDefSet loop
    go hoisted kept seen [] = (reverse hoisted, reverse kept)
    go hoisted kept seen (s@(MOVE (TEMP t) e) : rest)
      | canHoist t e seen =
          go (s : hoisted) kept (s : seen) rest
    go hoisted kept seen (s : rest) =
          go hoisted (s : kept) (s : seen) rest

    canHoist t e seen =
      isPure e
      && Set.null (freeTempsExpr e `Set.intersection` defs)
      && defCount t (loopBodyStmts loop) == 1
      && Set.notMember t (Set.unions (map stmtReadTemps seen))
\end{minted}

A função \texttt{spliceLoop} reconstrói a lista de instruções, inserindo os
comandos içados imediatamente antes do \texttt{LABEL lh}:

\begin{minted}{haskell}
spliceLoop stmts loop hoisted newBody =
  before ++ hoisted ++ prefix ++ newBody ++ suffix ++ after
  where
    before = take (loopHeadIdx loop) stmts
    ...
\end{minted}

\subsection{Desenrolamento: \texttt{IR.Opt.Loop.Unroll}}\label{sec:impl-unroll}

O desenrolamento usa uma função \texttt{makeCopy} para gerar cada cópia extra:

\begin{minted}{haskell}
makeCopy copyIdx =
  let suffix    = "_u" ++ show copyIdx
      renameMap = Map.fromList
        [ (l, l ++ suffix) | l <- intLabels ]
      entryLb   = lh ++ "_unroll_" ++ show copyIdx
      renamedBody = map (renameLabels renameMap) body
  in [CJUMP cond entryLb lf, LABEL entryLb] ++ renamedBody
\end{minted}

O laço resultante concatena o corpo original com as cópias extras:

\begin{minted}{haskell}
newLoopSlice =
  [ LABEL lh, CJUMP cond lb lf, LABEL lb ]
  ++ body
  ++ concatMap makeCopy [1 .. k - 1]
  ++ [ JUMP (NAME lh), LABEL lf ]
\end{minted}

\subsection{Fusão: \texttt{IR.Opt.Loop.Fusion}}\label{sec:impl-fusion}

A condição de fusibilidade é verificada pela função \texttt{areFusible}:

\begin{minted}{haskell}
areFusible l1 l2 =
  loopCond l1 == loopCond l2
  && not (null (loopBodyStmts l1))
  && not (null (loopBodyStmts l2))
  && last (loopBodyStmts l1) == last (loopBodyStmts l2)
  && noCrossDep l1 l2

noCrossDep l1 l2 =
  Set.null (defs1NoIndVar `Set.intersection` firstUses (loopBodyStmts l2))
  where
    indVars       = freeTempsExpr (loopCond l1)
    defs1NoIndVar = loopDefSet l1 `Set.difference` indVars
\end{minted}

A fusão propriamente dita em \texttt{fuseTwo} monta o corpo fundido removendo a
última instrução (atualização da variável de indução) do primeiro corpo:

\begin{minted}{haskell}
fuseTwo stmts l1 l2 = before ++ fusedSlice ++ after
  where
    body1NoUpdate = init (loopBodyStmts l1)
    body2         = loopBodyStmts l2
    fusedBody     = body1NoUpdate ++ body2
    fusedSlice    =
      [ LABEL lh, CJUMP cond lb lf, LABEL lb ]
      ++ fusedBody
      ++ [ JUMP (NAME lh), LABEL lf ]
\end{minted}

\subsection{Inlining: \texttt{IR.Opt.Inline}}\label{sec:impl-inline}

O módulo \texttt{IR.Opt.Inline} expõe \texttt{inlineProgram} e a configuração
\texttt{InlineConfig}, que controla o limiar de tamanho $\tau$:

\begin{minted}{haskell}
data InlineConfig = InlineConfig
  { inlineThreshold :: Int   -- número máximo de instruções linearizadas
  } deriving (Eq, Show)

defaultInlineConfig :: InlineConfig
defaultInlineConfig = InlineConfig { inlineThreshold = 10 }
\end{minted}

A entrada de \texttt{inlineProgram} determina quais funções são elegíveis e as
aplica em cada definição:

\begin{minted}{haskell}
inlineProgram :: InlineConfig -> Program -> Program
inlineProgram cfg prog =
  removeDeadFuncs $
    map (inlineFuncDef funcMap inlineable) prog
  where
    funcMap    = Map.fromList [(funcName fd, fd) | fd <- prog]
    inlineable = Set.fromList
      [ funcName fd
      | fd <- prog
      , funcName fd /= "main"
      , funcSize fd <= inlineThreshold cfg
      , not (isSelfRecursive fd) ]
\end{minted}

A função \texttt{isSelfRecursive} coleta todos os \texttt{NAME} usados em
posição de chamada e verifica se o nome da função está entre eles.
\texttt{removeDeadFuncs} filtra definições que não são mais referenciadas por
nenhuma outra função (exceto \texttt{main}).

A reescrita propriamente dita usa uma mônada de estado para gerar sufixos de
frescura únicos (\texttt{InlineM = State Int}):

\begin{minted}{haskell}
inlineValueCall :: Map Label FuncDef -> Set Label
                -> Label -> [Expr] -> InlineM Expr
inlineValueCall fm il f args = do
  n <- freshId
  let suf     = "_il" ++ show n
      fd      = fm Map.! f
      locals  = localLabels (funcBody fd)
      body0   = renameLocals suf locals (funcBody fd)
      retTemp = "_il_ret" ++ suf
      exitLbl = "L_il_exit" ++ suf
      body1   = transformReturns retTemp exitLbl body0
      il'     = Set.delete f il   -- evita re-inlining do mesmo f
  args'  <- mapM (rewriteExpr fm il) args
  body2  <- rewriteStmt fm il'     body1
  let pMoves = paramMoveList (map (++ suf) (funcParams fd)) args'
      full   = seqStmts (pMoves ++ [body2, LABEL exitLbl])
  return (ESEQ full (TEMP retTemp))
\end{minted}

A função \texttt{renameLocals} percorre toda a árvore de comandos e expressões,
adicionando o sufixo $s$ a cada temporário e a cada rótulo que seja \emph{local}
à função (definido por \texttt{LABEL} dentro de seu corpo).  Rótulos externos
(nomes de outras funções em chamadas) são preservados intactos.
\texttt{transformReturns} percorre o corpo e substitui cada
$\mathbf{RETURN}([e])$ pelo par $\mathbf{MOVE}(r^s, e) \mathbin{;}
\mathbf{JUMP}(\ell_\mathit{exit}^s)$.

Para chamadas void (em posição de comando), \texttt{inlineVoidCall} produz um
\texttt{Stmt} diretamente, evitando o envoltório \texttt{ESEQ}.

\subsection{Eliminação de Chamadas em Cauda:
\texttt{IR.Opt.TCO}}\label{sec:impl-tco}

O módulo \texttt{IR.Opt.TCO} implementa a Definição~\ref{def:tco} em duas
etapas: detecção e transformação.

A detecção procura os dois padrões de cauda na lista linearizada:

\begin{minted}{haskell}
anyTailCall :: Label -> [Stmt] -> Bool
anyTailCall f stmts =
  any isTCO stmts ||
  any isVoidTCO (zip stmts (drop 1 stmts))
  where
    isTCO (RETURN [CALL (NAME g) _])          = f == g
    isTCO _                                    = False
    isVoidTCO (EXP (CALL (NAME g) _), RETURN []) = f == g
    isVoidTCO _                                   = False
\end{minted}

A transformação gera os temporários auxiliares e reescreve os padrões
identificados:

\begin{minted}{haskell}
tcoFuncDef :: FuncDef -> FuncDef
tcoFuncDef fd
  | not (anyTailCall (funcName fd) stmts) = fd
  | otherwise = fd { funcBody = rebuildSeq (LABEL lh : transformed) }
  where
    stmts       = linearize (funcBody fd)
    lh          = "L_tco_" ++ funcName fd
    transformed = transformStmts (funcName fd) (funcParams fd) lh stmts

transformStmts :: Label -> [Temp] -> Label -> [Stmt] -> [Stmt]
transformStmts f params lh = go
  where
    n        = length params
    tcoTemps = ["_tco_" ++ show i | i <- [0 .. n - 1]]

    makeTCO :: [Expr] -> [Stmt]
    makeTCO args =
      zipWith (\t e -> MOVE (TEMP t) e) tcoTemps args ++
      zipWith (\p t -> MOVE (TEMP p) (TEMP t)) params tcoTemps ++
      [JUMP (NAME lh)]

    go [] = []
    go (RETURN [CALL (NAME g) args] : rest)
      | g == f = makeTCO args ++ go rest
    go (EXP (CALL (NAME g) args) : RETURN [] : rest)
      | g == f = makeTCO args ++ go rest
    go (s : rest) = s : go rest
\end{minted}

Os temporários auxiliares $\mathtt{\_tco\_0}, \ldots$ implementam o passo duplo
da Definição~\ref{def:tco}, quebrando dependências cíclicas entre parâmetros.
A função percorre a lista de comandos linearizados com \texttt{go}, consumindo
dois elementos de uma vez no padrão void (para que o \texttt{RETURN []} não seja
emitido separadamente).

\subsection{Pipeline de Otimizações:
\texttt{IR.Opt.Pipeline}}\label{sec:impl-pipeline}

O módulo \texttt{IR.Opt.Pipeline} compõe todas as otimizações e define o
pipeline completo.  O utilitário \texttt{fixedPoint} do módulo
\texttt{Utils.Fixpoint} aplica uma transformação até que o resultado se
estabilize:

\begin{minted}{haskell}
optAll :: Program -> Program
optAll =
    dceProgram
  . cseProgram
  . unrollProgram defaultUnrollConfig
  . fixedPoint fuseProgram
  . cseProgram
  . fixedPoint foldPropProgram
  . licmProgram
  . fixedPoint foldPropProgram
  . fixedPoint (inlineProgram defaultInlineConfig)
  . tcoProgram
\end{minted}

A composição respeita as dependências entre otimizações:
\begin{enumerate}
  \item \texttt{tcoProgram} converte auto-recursão em cauda em laços, expondo
    estrutura iterativa para as otimizações seguintes.
  \item \texttt{inlineProgram} a ponto fixo substitui chamadas pequenas pelos
    respectivos corpos; funções mortas são removidas.  Pode expor novas
    constantes e novas chamadas em cauda.
  \item \texttt{foldPropProgram} a ponto fixo aproveita as constantes expostas
    pelo inlining.
  \item \texttt{licmProgram} move código invariante para fora dos laços criados
    pela TCO e pelo inlining.
  \item Uma segunda rodada de \texttt{foldPropProgram} dobra constantes
    içadas.
  \item \texttt{cseProgram} elimina subexpressões duplicadas.
  \item \texttt{fuseProgram} a ponto fixo combina laços adjacentes.
  \item \texttt{unrollProgram} replica corpos de laços.
  \item \texttt{cseProgram} e \texttt{dceProgram} limpam o código após o
    desenrolamento.
\end{enumerate}

Cada otimização isolada também pode ser invocada pelo executável \texttt{ir}:

\begin{minted}{text}
ir --opt-fold      arquivo.ir   # dobramento de constantes
ir --opt-prop      arquivo.ir   # propagação + dobramento
ir --opt-cse       arquivo.ir   # eliminação de subexpressões comuns
ir --opt-dce       arquivo.ir   # eliminação de código morto
ir --opt-licm      arquivo.ir   # movimentação de invariante de laço
ir --opt-unroll 4  arquivo.ir   # desenrolamento com fator 4
ir --opt-fusion    arquivo.ir   # fusão de laços
ir --opt-inline    arquivo.ir   # inlining de funções
ir --opt-tco       arquivo.ir   # eliminação de chamadas em cauda
ir --opt-all       arquivo.ir   # pipeline completo
\end{minted}

% ============================================================================
\section{Conclusão}\label{sec:opt-conclusao}
% ============================================================================

Este capítulo apresentou formalmente dez otimizações aplicadas à representação
intermediária em árvore (IRT): dobramento de constantes, propagação de
constantes, eliminação de subexpressões comuns, eliminação de código morto,
movimentação de código invariante de laço, desenrolamento, fusão de laços,
inlining de funções e eliminação de chamadas em cauda.

Para cada otimização foram apresentados:
\begin{itemize}
  \item a definição formal em termos da IRT;
  \item as condições de correção que garantem a equivalência observacional
    entre o programa original e o otimizado;
  \item exemplos concretos na IRT;
  \item a implementação em Haskell, destacando a correspondência entre
    teoria e código.
\end{itemize}

As otimizações são compostas em um pipeline que aplica cada transformação na
ordem que maximiza as oportunidades de melhoria.  O uso de ponto fixo garante
que otimizações mutuamente dependentes (como dobramento e propagação) são
aplicadas até a convergência.
